\chapter{Эффективное действие изогнутой вихревой нити}
\label{ch:v6-bosonic-defect-curved-string-effective-action-gate}

\section{Постановка}

Полный поперечный оператор прямой нити устойчив и имеет внутреннюю щель
$3.61933633$. Следующий вопрос --- превращается ли такая нить в
локализованный объект после замыкания в кольцо. Для этого недостаточно
устойчивости поперечного сечения: требуется вывести действие мировой
поверхности и проверить энергию как функцию радиуса кольца.

Систематические выводы эффективных действий вихрей различают переносные
голдстоуновские моды и массивные деформации профиля. Первая подсистема даёт
действие Намбу--Гото; нетривиальный кривизный ответ требует интегрирования
массивных мод \cite{curved-string-aurrekoetxea,curved-string-anderson}.
Более ранние коллективно-координатные подходы также находили зависимость от
внешней кривизны, но её знак и порядок зависят от режима и способа устранения
тяжёлых полей \cite{curved-string-orland,curved-string-eto}.

\section{Поперечный энергетический масштаб}

Из стационарного четырёхпрофильного решения повторно интегрирована плотность
энергии. Получено
\begin{equation}
 T=1.5744530783,
\end{equation}
что совпадает с прежним натяжением с относительным остатком
$2.46\cdot10^{-7}$.

Радиальные моменты равны
\begin{equation}
 I_2=0.5174619005,
 \qquad
 I_4=0.4366199900,
 \qquad
 I_6=0.7206475643.
\end{equation}
Энергетическая среднеквадратичная ширина нити поэтому равна
\begin{equation}
 w=\sqrt{I_2/T}=0.5732899499.
\end{equation}
В радиусах $0.88022$, $1.04103$ и $1.39991$ сосредоточены соответственно
$90\%$, $95\%$ и $99\%$ энергии поперечного сечения.

\section{Нулевая мода и действие Намбу--Гото}

При медленном переносе центра профиль зависит от двух нормальных координат
относительно мировой поверхности. Осесимметрия обращает первый поперечный
момент энергии в нуль. Поэтому ведущая редукция имеет вид
\begin{equation}
 S_0=-T\int d^2\sigma\,\sqrt{-\gamma}.
\end{equation}

Самостоятельный квадратичный член внешней жёсткости из одних переносных мод
не возникает. Возможный член, пропорциональный внутренней скалярной кривизне
мировой поверхности, топологичен; для статического цилиндра
$\mathbb R\times S^1$ его локальная кривизна равна нулю.

Проект также не содержит нормируемой внутренней окружности или дионной фазы:
ориентационный ротор расходится, а локальное пространство модулей прямой
нити состоит только из переносов. Следовательно, на мировой поверхности нет
выведенного постоянного тока, который мог бы дать в энергии кольца член
$1/R$.

\section{Энергия спокойного кольца}

Для статического кольца радиуса $R$ действие Намбу--Гото даёт
\begin{equation}
 E_0(R)=2\pi T R,
 \qquad
 \frac{dE_0}{dR}=2\pi T=9.89258045>0.
\end{equation}
Следовательно, ведущая энергия не имеет стационарного радиуса: кольцо
стремится сжаться.

При $R/w=3,5,10$ получены
\begin{equation}
 E_0=17.01395,
 \qquad 28.35658,
 \qquad 56.71317.
\end{equation}
Эти числа не являются массами частиц: общий размерный масштаб модели ещё не
восстановлен.

\section{Граница поправок конечной толщины}

В общем виде первый нетривиальный динамический ответ может начинаться как
\begin{equation}
 E(R)=2\pi T R
 \left[1+c_4\left(\frac wR\right)^4+\cdots\right].
\end{equation}
Для $c_4=1$ формальный стационарный радиус равен
\begin{equation}
 \frac{R_*}{w}=3^{1/4}=1.31607,
\end{equation}
то есть лежит вне режима медленной кривизны $R\gg w$. Чтобы получить
стационарную точку уже при $R=3w$, потребовался бы коэффициент
$c_4=27$; при $R=5w$ --- $c_4=208.33$.

Ни натяжение, ни внутренняя щель сами по себе не определяют $c_4$. Для его
вывода требуется решить неоднородные уравнения массивных поперечных мод на
изогнутом фоне.

\begin{theorem}[No-go спокойного кольца в нулемодовом приближении]
Текущий родитель строго выводит ведущую энергию спокойной замкнутой нити
$E_0(R)=2\pi TR$, но не выводит стабилизирующий член порядка $1/R$.
Следовательно, в контролируемом низкоэнергетическом приближении основное
состояние прямой нити нельзя повысить до устойчивой частицы-кольца: оно
сжимается под действием натяжения.
\end{theorem}

\begin{nogocriterion}
No-go не исключает хопфово скрученную или возбуждённую петлю. Он исключает
только спокойное кольцо, построенное из переносных нулевых мод без
продольного заряда и без вычисленного ответа массивных полей. Объявлять
частицу до расчёта этих каналов нельзя.
\end{nogocriterion}

\section{Следующий различающий тест}

Нулевая гипотеза: кривизна источает критическую массивную пару со щелью
$3.61933633$; решение неоднородного поперечного уравнения фиксирует знак и
величину первого нетривиального коэффициента $c_4$ без нового параметра.

Стоп-критерий: если полученный коэффициент не создаёт контролируемого
минимума при $R\gtrsim3w$, маршрут спокойной замкнутой петли закрывается.
Тогда для частицы потребуется отдельный топологический поворот, продольное
возбуждение или иной уже выведенный сохраняющийся заряд.

\begin{protocol}
\begin{enumerate}
 \item натяжение прямой нити: \textbf{$1.57445308$};
 \item энергетическая ширина: \textbf{$0.57328995$};
 \item ведущая мировая поверхность: \textbf{Намбу--Гото};
 \item независимая квадратичная жёсткость нулевых мод: \textbf{нет};
 \item нормируемый внутренний ток: \textbf{нет};
 \item спокойное кольцо при ведущем порядке: \textbf{сжимается};
 \item коэффициент $c_4$: \textbf{не вычислен};
 \item контролируемый конечный радиус: \textbf{не получен};
 \item хопфово скрученная петля: \textbf{не исключена};
 \item рождение материи: \textbf{не закрыто}.
\end{enumerate}
\end{protocol}

Машинный сертификат сохранён в
\path{s2t_v6_bosonic_defect_curved_string_effective_action_gate_results.json}.

\begin{thebibliography}{9}
\bibitem{curved-string-aurrekoetxea}
J.~C.~Aurrekoetxea, J.~J.~Blanco-Pillado,
A.~Garc\'ia Mart\'in-Caro, J.~M.~Queiruga,
\emph{On curvature corrections for field theory cosmic strings},
arXiv:2603.05243.

\bibitem{curved-string-anderson}
M.~Anderson, F.~Bonjour, R.~Gregory, J.~Stewart,
\emph{Effective action and motion of a cosmic string},
arXiv:hep-ph/9707324.

\bibitem{curved-string-orland}
P.~Orland,
\emph{Extrinsic Curvature Dependence of Nielsen--Olesen Strings},
arXiv:hep-th/9404140.

\bibitem{curved-string-eto}
M.~Eto, T.~Fujimori, M.~Nitta, K.~Ohashi, N.~Sakai,
\emph{Higher Derivative Corrections to Non-Abelian Vortex Effective Theory},
arXiv:1204.0773.
\end{thebibliography}